% Fragmento público generado desde propietarios íntegros.
% Residencia: 52.3.2.
% Fuente: ../incoming/radion_monografia_20260827/manuscrito/sections/04_cincuenta_critica.tex:106-171
\subsubsection{Typed separation of the 2 objects with value \(1/2\)}

APP contains, in addition, a singular mode with
\[
s=\frac12,\qquad s^2=\frac14,\qquad
1-s^2=\frac34=\frac{90}{120},\qquad
s\sqrt{1-s^2}=\frac{\sqrt3}{4}.
\]
This \(s\) arises from the principal angles of the sum and product sheets. The
critical seam \(\Re\rho=1/2\) arises from a functional involution. That both
coordinates equal \(1/2\) is a scalar coincidence with possible content,
but it does not authorize their identification without a typed square.

The correct way to preserve the finding is an equalizer diagram:
\[
\begin{tikzpicture}[baseline=(current bounding box.center),node distance=14mm and 23mm,
 every node/.style={font=\small}]
\node[draw=RadBlue,rounded corners,fill=RadMist] (crit) {critical seam \(\Re s=1/2\)};
\node[draw=RadTeal,rounded corners,fill=RadCream,right=of crit] (face) {chart \(j_{100}\)};
\node[draw=RadCopper,rounded corners,fill=RadCream,right=of face] (theta) {phase \(\theta_2=50^\circ\)};
\node[draw=RadRose,rounded corners,fill=white,below=of face] (mode) {APP mode \(s=1/2\), \(1-s^2=90/120\)};
\draw[-{Latex},thick,RadBlue] (crit)--node[above]{\(j_{100}\)}(face);
\draw[-{Latex},thick,RadTeal] (face)--node[above]{equalizer}(theta);
\draw[-{Latex},dashed,RadRose] (mode)--node[right,align=left]{same figure;\\distinct domain}(face);
\end{tikzpicture}
\]
The diagram does not feign the missing arrow. It shows what has been proved, what has
been formalized, and what equality still requires additional transport if
a global operator identification is intended.

\subsubsection{Spectral Circle and Helical Memory Compatibility}

The Cayley--Pontryagin reader transports a spectral fiber via
\[
\widetilde U_{\Gamma_9^n}J\psi_\gamma
=\tau_\gamma^nJ\psi_\gamma.
\]
The circular publication is \(\tau_\gamma^n\in S^1\). The lift
\[
(\tau_\gamma^n,n)\in S^1\times\Z
\]
preserves the number of returns. On the contractive sheet appears
\[
\widetilde M_9^nJ\psi_\gamma
=\left(\frac59\right)^n\tau_\gamma^nJ\psi_\gamma.
\]
We therefore have three typed images:
\[
\text{phase circle},\qquad
\text{memory helix},\qquad
\text{contractive spiral}.
\]

A zero is not a revolution, and a nine is not a particular zero of the zeta
function. The zero fixes a character \(\tau_\gamma\) that persists through the
revolutions; nonadic return increments memory. Nine functions as residual zero
in \(\Z/9\Z\), an arithmetic operation distinct from the enumeration
of spectral zeros.

\begin{narrativebox}
The critical line becomes a circle when we forget the linear height and
retain the unitary character. It becomes a helix when we reinsert
the memory of how many times the character has been transported. Depth
is not a third decorative coordinate: it is the information that the circular
projection cannot retain.
\end{narrativebox}
